\section{\texorpdfstring{$SU(3)$ 与夸克模型}{SU(3) 与夸克模型}}
\label{chap:su3-quarks}

最高权把不可约表示的分类化成了一组离散标签；要把这些标签用于计算，还要能够画出权图、求维数并分解张量积。秩为二的 $SU(3)$ 是第一个能在平面上看见完整根系与权格的例子，同时又直接作用在强子谱和夸克颜色指标上。

两个单根张成 \(SU(3)\) 的根平面，两个基本权则标记全部有限维不可约表示。最低的几个权图已经足以分解 \(\mathbf3\otimes\mathbf3\) 与 \(\mathbf3\otimes\bar{\mathbf3}\)，并解释夸克怎样组成介子和重子。味 \(SU(3)\) 是近似的全局对称，色 \(SU(3)_c\) 是局域规范对称；群结构相同，并不意味着它们作用在同一物理空间。

\subsection{\texorpdfstring{$SU(3)$ 的基本表示}{SU(3) 的基本表示}}

$SU(3)$ 的秩为 $2$，不可约表示可用 Dynkin 标号 $(p,q)$ 标记。取归一化 $(\alpha_i,\alpha_i)=1$、$\langle\alpha_1,\alpha_2\rangle=120^\circ$，则两个单根在平面中具体为
\[
\alpha_1=(1,0),\qquad \alpha_2=\left(-\tfrac12,\tfrac{\sqrt3}{2}\right),
\]
第三个正根 $\alpha_1+\alpha_2=(\tfrac12,\tfrac{\sqrt3}{2})$。按定义 $2(\omega_i,\alpha_j)/(\alpha_j,\alpha_j)=\delta_{ij}$ 解得基本权
\[
\omega_1=\left(\tfrac12,\tfrac{\sqrt3}{6}\right),\qquad
\omega_2=\left(0,\tfrac{\sqrt3}{3}\right),
\]
且 $\rho=\omega_1+\omega_2=\alpha_1+\alpha_2$。把 $\pm\alpha_1,\pm\alpha_2,\pm(\alpha_1+\alpha_2)$ 六个根画在平面上构成一个六边形：$\alpha_1$ 指向正东，$\alpha_2$ 指向西北偏上，相邻根夹角 $120^\circ$，对径根相差一个负号；基本权 $\omega_1,\omega_2$ 分别落在两个 Weyl 室内，与相应单根成 $30^\circ$。不可约表示的维数为
\[
\dim(p,q)=\frac12(p+1)(q+1)(p+q+2).
\]
基本表示 $\mathbf3=(1,0)$，共轭表示 $\bar{\mathbf3}=(0,1)$，伴随表示 $\mathbf8=(1,1)$。夸克味三重态可写作 $q=(u,d,s)^{\mathsf T}$；反夸克处在 $\bar{\mathbf3}$ 中。介子和重子的母分解分别为
\[
\mathbf3\otimes\bar{\mathbf3}=\mathbf8\oplus\mathbf1,
\]
以及
\begin{equation}
\mathbf3\otimes\mathbf3\otimes\mathbf3
=\mathbf{10}\oplus\mathbf8\oplus\mathbf8\oplus\mathbf1.
\label{eq:baryon-decomposition}
\end{equation}
等式左边的任意元素可写成矩阵 $M^i{}_j=q^i\bar q_j$。按迹分解
\[
M^i{}_j=\left(M^i{}_j-\frac13\delta^i{}_j M^k{}_k\right)
+\frac13\delta^i{}_j M^k{}_k .
\]
在 $M\mapsto U M U^\dagger$ 下，第二项保持不变，是一维平凡表示；第一项无迹，有 $3^2-1=8$ 个独立分量，并按伴随表示变换。因此
\(\mathbf3\otimes\bar{\mathbf3}=\mathbf8\oplus\mathbf1\) 不是单凭维数猜出的分解，而是由不变迹映射及其核直接得到。

第一步 $\mathbf3\otimes\mathbf3=\mathbf6_s\oplus\bar{\mathbf3}_a$ 区分前两个指标的对称性；再乘第三个 $\mathbf3$，分别得到 $\mathbf{10}\oplus\mathbf8$ 与 $\mathbf8\oplus\mathbf1$。Young 图同时记录这一置换对称性，因此分解不需要逐个猜测粒子态。

\begin{theorem}[Gell--Mann--Okubo 质量公式]
若 $SU(3)$ 味对称只被夸克质量的一阶八重态扰动破坏，八重态重子质量满足
\begin{equation}
\frac{M_N+M_\Xi}{2}=\frac{3M_\Lambda+M_\Sigma}{4},
\label{eq:gmo-octet}
\end{equation}
十重态则给出等间距关系
\begin{equation}
M_{\Sigma^*}-M_\Delta
=M_{\Xi^*}-M_{\Sigma^*}
=M_\Omega-M_{\Xi^*}.
\label{eq:gmo-decuplet}
\end{equation}
\end{theorem}

\begin{proof}
八重态扰动在一个给定八重态表示中的对角矩阵元，只能由等标量
$Y$ 与 $I(I+1)-Y^2/4$ 的两个独立一阶结构组成。因此可写
\begin{equation}
M(Y,I)=M_0+aY+b\left[I(I+1)-\frac{Y^2}{4}\right].
\label{eq:gmo-mass-parametrization}
\end{equation}
对核子 $N$、$\Lambda$、$\Sigma$、$\Xi$，量子数分别为
\[
(Y,I)=
\left(1,\frac12\right),(0,0),(0,1),
\left(-1,\frac12\right).
\]
代入\eqnrefs{eq:gmo-mass-parametrization} 得
\[
M_N=M_0+a+\frac b2,\quad
M_\Lambda=M_0,\quad
M_\Sigma=M_0+2b,\quad
M_\Xi=M_0-a+\frac b2.
\]
消去 $M_0,a,b$，
\[
\frac{M_N+M_\Xi}{2}
=M_0+\frac b2
=\frac{3M_\Lambda+M_\Sigma}{4}.
\]

十重态中每个超荷只出现一个等同位旋多重态，一阶八重态扰动的
对角矩阵元因而是 $Y$ 的仿射函数。相邻成员
$\Delta,\Sigma^*,\Xi^*,\Omega$ 的超荷依次相差 $1$，质量差便相等。
二阶破缺允许更多独立群论结构，上述线性关系也随之得到修正。
\end{proof}

\begin{proof}[$\mathbf3^{\otimes3}$ 的分解]
Young 图的形状用行长列表 $[\lambda_1,\lambda_2,\ldots]$ 表示。添盒规则给出
\[
[1]\otimes[1]=[2]\oplus[1,1],
\]
对应 $\mathbf6_s\oplus\bar{\mathbf3}_a$。再添一个盒子，
\[
[2]\otimes[1]=[3]\oplus[2,1],
\qquad
[1,1]\otimes[1]=[2,1]\oplus[1,1,1].
\]
四个图的 $SU(3)$ 维数依次为 $10,8,8,1$，总维数 $27=3^3$，得到\eqnrefs{eq:baryon-decomposition}。两个八重态同构，但来自不同的指标置换对称性。
\end{proof}

\subsection{颜色与味}

同一个抽象群 $SU(3)$ 可以作用在不同向量空间上。味变换混合
$u,d,s$ 三种轻夸克，近似性来自它们质量不完全相等；颜色变换混合
$r,g,b$ 三个颜色分量，是规范对称。一个物理强子可以属于非平凡的
味多重态，却必须在颜色空间中成为单态。

先看介子。若夸克颜色为 $q^a$、反夸克颜色为 $\bar q_b$，则
\[
\mathbf3_c\otimes\bar{\mathbf3}_c
=\mathbf1_c\oplus\mathbf8_c.
\]
归一化颜色单态为
\[
|1_c\rangle=\frac1{\sqrt3}
\left(|r\bar r\rangle+|g\bar g\rangle+|b\bar b\rangle\right).
\]
在 $U\in SU(3)_c$ 下，
\[
\delta^a{}_b q^b\bar q_a
\longmapsto
\delta^a{}_b U^b{}_c q^c
\bar q_d(U^\dagger)^d{}_a
=\delta^c{}_d q^d\bar q_c.
\]
所以不变张量 $\delta^a{}_b$ 把
$\mathbf3_c\otimes\bar{\mathbf3}_c$ 投影到单态。

三个夸克的颜色单态由另一个不变张量给出：
\begin{equation}
|1_c\rangle_{\mathrm{baryon}}
=\frac1{\sqrt6}\epsilon_{abc}|a\,b\,c\rangle.
\label{eq:baryon-color-singlet}
\end{equation}
因为
\[
\epsilon_{abc}U^a{}_{a'}U^b{}_{b'}U^c{}_{c'}
=(\det U)\epsilon_{a'b'c'}
=\epsilon_{a'b'c'},
\]
\eqnrefs{eq:baryon-color-singlet} 在颜色变换下不变。它在任意两夸克
交换下变号。对三个处于同一空间基态的夸克，总波函数必须反对称；
颜色部分已经反对称，因此自旋与味的联合部分必须对称。
$\Delta^{++}=uuu$ 的味部分完全对称，自旋部分便只能取对称的
$j=3/2$ 多重态。
